\documentclass[compress]{beamer}
\usetheme{Warsaw}
\usecolortheme{crane}
\usepackage[utf8]{inputenc}
\usepackage[portuguese]{babel}
\usepackage{url}
\usepackage{tikz}
\usepackage{color}
\usetikzlibrary{arrows}
\usetikzlibrary{arrows}
\usepackage{default}
\usepackage{amsmath}
\usepackage{amsfonts}
\usepackage{amssymb}
\usepackage{newalg}

% Algumas definições %%%
\def\term#1{{\sc #1}}   % IR terms in examples not index terms!
\def\query#1{{\sf #1}}
\def\oper#1{{\sc #1}} % AND, OR, NOT
%%%%%%%%%%%%%%%%%%%%%%%%

\title[Coelho: Vocabulário de termos e lista de ``postings'']
{Introdução à Recuperação de Informações\\
\large \url{https://github.com/fccoelho/curso-IRI}\\[0.5cm]
IRI 2: Vocabulário de termos e lista de ``postings''}

\author [Coelho F.C. \& Souza R.R.]{ Flávio Codeço Coelho}

\institute [EMAp, FGV]{Escola de Matemática Aplicada,   Fundação Getúlio Vargas}
\date


\begin{document}

\begin{frame}
\titlepage
\end{frame}

\begin{frame}[fragile]
\frametitle{Sumário da Aula}
\tableofcontents
\end{frame}
\section{Recapitulação}
\begin{frame}
\frametitle{Índice invertido}
Para cada termo $t$, armazenamos uma lista de todos os documentos que contém 
$t$.

\bigskip
\begin{tabular}{|c|c|r|r|r|r|r|r|r|r|r|}
\cline{1-1}\cline{3-10}
\term{Brutus} & $\longrightarrow$ & 1 & 2 & 4 & 11 & 31 & 45 & 173 & 174 \\ 
\cline{1-1}\cline{3-10}
\multicolumn{8}{l}{} \\ \cline{1-1}\cline{3-11}
\term{Caesar} & $\longrightarrow$ & 1 & 2 & 4 & 5 & 6 & 16 & 57 & 132 & \ldots 
\\ \cline{1-1}\cline{3-11}
\multicolumn{8}{l}{} \\ \cline{1-1}\cline{3-6}
\term{Calpurnia} & $\longrightarrow$ & 2 & 31 & 54 & 101 \\
\cline{1-1}\cline{3-6} \multicolumn{8}{l}{}  \\
% The multicolumn{1}'s below suppress the vertical lines....
\multicolumn{1}{c}{$\vdots$} \\
\multicolumn{1}{c}{$\underbrace{\phantom{\mbox{Calpurnia}}}$} &
\multicolumn{1}{c}{} &
\multicolumn{9}{c}{$\underbrace{\phantom{\mbox{Calpurnia Calpurnia
Calpurnia Caesar hath}}}$} \\
\multicolumn{1}{c}{\visible<1>{\textbf{dictionary}}} &
\multicolumn{1}{c}{} & \multicolumn{9}{c}{\visible<1>{\textbf{postings}}}
\end{tabular}
\end{frame}

\begin{frame}
\frametitle{Interseção de duas listas de ``postings''}

\begin{tabular}{lll}
\term{Brutus} & $\longrightarrow$ &
\alert<2>{\framebox{1}}$\rightarrow$\alert<3,4>{\framebox{2}}
$\rightarrow$\alert<5>{\framebox{4}}$\rightarrow$\alert<6>{\framebox{11}}
$\rightarrow$\alert<7,8>{\framebox{31}}$\rightarrow$\alert<9>{\framebox{45}}
$\rightarrow$\alert<10,11>{\framebox{173}}$\rightarrow$\framebox{174}\\[1ex]
\term{Calpurnia} & $\longrightarrow$ &
\alert<2-4>{\framebox{2}}$\rightarrow$\alert<5-8>{\framebox{31}}
$\rightarrow$\alert<9-10>{\framebox{54}}$\rightarrow$\alert<11>{\framebox{101}}
\\[2ex]
Intersection & $\Longrightarrow$ & 
\visible<4->{\framebox{2}}\visible<8->{$\rightarrow$\framebox{31}}
\end{tabular}

\end{frame}

\begin{frame}
\frametitle{Construindo o índice invertido: ordenando os ``postings''}
\begin{tiny}
\mbox{
\begin{tabular}{@{}lr}
\textbf{term} & \textbf{\llap{doc}ID} \\
I & 1 \\
did & 1 \\
enact & 1 \\
julius & 1 \\
caesar & 1 \\
I & 1 \\
was & 1 \\
killed & 1 \\
i' & 1 \\
the & 1 \\
capitol & 1 \\
brutus & 1 \\
killed & 1 \\
me & 1 \\
so & 2 \\
let & 2 \\
it & 2 \\
be & 2 \\
with & 2 \\
caesar & 2 \\
the & 2 \\
noble & 2 \\
brutus & 2 \\
hath & 2 \\
told & 2 \\
you & 2 \\
caesar & 2 \\
was & 2 \\
ambitious & 2
\end{tabular}
%
{\huge $\Longrightarrow$}
%
\begin{tabular}{lr}
\textbf{term} & \textbf{\llap{doc}ID} \\
ambitious & 2 \\
be & 2 \\
brutus & 1 \\
brutus & 2 \\
capitol & 1 \\
caesar & 1 \\
caesar & 2 \\
caesar & 2 \\
did & 1 \\
enact & 1 \\
hath & 1 \\
I & 1 \\
I & 1 \\
i' & 1 \\
it & 2 \\
julius & 1 \\
killed & 1 \\
killed & 1 \\
let & 2 \\
me & 1 \\
noble & 2 \\
so & 2 \\
the & 1 \\
the & 2 \\
told & 2 \\
you & 2 \\
was & 1 \\
was & 2 \\
with & 2
\end{tabular}
}
\end{tiny}
\end{frame}

\begin{frame}
\frametitle{O Google usa busca Booleana?}
\begin{itemize}
\item No Google, a interpretação defaul de uma consulta 
  [$w_1$ $w_2$ 
\ldots $w_n$] é 
  $w_1$ AND $w_2$ AND 
\ldots AND $w_n$
\item Casos em que você recebe resultados que não contém uma das $w_i$:
\begin{itemize}
\item âncora de texto
\item página contém variantes da $w_i$ (morfologia, correção ortográfica
 , sinônimo)
\item consultas longas ($n$ é grande)
\item expressões booleanas geram poucos resultados.
\end{itemize}

\item Booleana simples vs.\ Ordenação do conjunto de resultados

\begin{itemize}
\item A busca booleana simples não gera ordenação dos documentos.
\item O Google (e a maioria dos buscadores booleanos bem feitos) ordena o 
conjunto de resultados -- com os melhores resultados (de acordo com um 
estimador de relevância) no topo.

\end{itemize}
\end{itemize}
\end{frame}

\section{Documentos}
\begin{frame}
\frametitle{Documentos}
\begin{itemize}[<+->]
\item Última aula: Sistema de recuperação Booleana simples
\item Nossos pressupostos eram:
\begin{itemize}[<+->]
\item Nós sabemos o que é um documento.
\item Documentos são ``legíveis por máquina''.
\end{itemize}
\item Na prática isto pode ser bem complicado.
\end{itemize}
\end{frame}

\begin{frame}
\frametitle{Parseando um Documento}
\begin{itemize}[<+->]
\item Precisamos lidar com formato e lingua de cada documento.
\item Em que formato ele está? pdf, word, excel, html etc.
\item Em que língua ele está?
\item Em que codificação está?
\item Cada uma destas perguntas implica um problema de 
classificação.
\item Alternativa: usar heurísticas.
\end{itemize}
\end{frame}

\begin{frame}
\frametitle{Complicações de formato e Língua}
\begin{itemize}[<+->]
\item Um mesmo índice normalmente contém termos de diversas línguas.
\begin{itemize}[<+->]
\item Algumas vezes um documento ou seus componentes contém múltiplas 
línguas e formatos.
\item Email em francês com PDF em espanhol anexado
\end{itemize}
\item Qual a unidade do documento para fins de indexação?
\item Um arquivo?
\item Um email?
\item Um email com 5 anexos?
\item Um grupo de arquivos(ppt ou latex in HTML)?
\item Em suma: Responder à pergunta ``O que é um documento?''
  não é trivial e requer algumas decisões.
\end{itemize}
\end{frame}

\section{Termos}

\subsection{Genericos + outras línguas}

\begin{frame}
\frametitle{Definições}
\begin{itemize}[<+->]
\item {\color{blue}Palavra} -- Uma string delimitada de caracteres como aparece 
no texto.
\item {\color{blue}Termo} -- Uma palavra ``normalizada'' (capitalização, 
morfologia, ortografia, etc); Uma classe equivalente de palavras.
\item {\color{blue}Token} -- Uma instancia de uma palavra ou termo ocorrendo em 
um documento.
\item {\color{blue}Tipo} -- O mesmo que termo na maioria dos casos: Uma classe 
de equivalência de tokens.
\end{itemize}
\end{frame}

\begin{frame}
\frametitle{Normalização}
\begin{itemize}[<+->]
\item Precisamos ``normalizar'' os termos no texto indexado, assim como o 
termos da consulta da mesma maneira.
\item Exemplo: Queremos associar \emph{U.S.A.} com \emph{USA}
\item A maneira mais comum é criar \myblue{classes de equivalência}
  de termos.
\item Alternativamente: fazer expansões assimétricas
\begin{itemize}[<+->]
\item window $\rightarrow$ window, windows
\item windows $\rightarrow$ Windows, windows
\item Windows (sem expansão)
\end{itemize}
\item Mais poderoso, porém menos eficiente
\item \mygreen{Porque não colocamos 
\emph{window},
\emph{Window},
\emph{windows}, e
\emph{Windows} na mesma classe de equivalência?}
\end{itemize}
\end{frame}

\begin{frame}
\frametitle{Normalização: outras línguas}
\begin{itemize}[<+->]
\item Normalização e detecção de linguagem interagem.
\item \emph{PETER WILL NICHT MIT.} $\rightarrow$ MIT
  $=$ mit
\item \emph{He got his PhD from MIT.} $\rightarrow$ MIT
  $\neq$ mit
\end{itemize}
\end{frame}

\begin{frame}
\frametitle{Lembre-se: Construção de índices invertidos}
\begin{itemize}[<+->]
\item Entrada:\\[0.2ex]
\framebox{\weestrut Friends, Romans, countrymen.}
\framebox{\weestrut So let it be with Caesar} \ldots
\item Saída:\\[0.2ex]
\framebox{\weestrut friend} \framebox{\weestrut roman}
  \framebox{\weestrut countryman} \framebox{\weestrut so} \ldots
\item Cada token é candidato a uma entrada de posting.
\item Quais são os tokens válidos?
\end{itemize}
\end{frame}

\begin{frame}
\frametitle{Exercícios}

\emph{In June, the dog likes to chase the cat in the barn.} --
{\color{green}Quantos tokens? Quantos tipos de Palavras?}

\bigskip

Porque a tokenização é difícil -- mesmo em inglês. {\color{green}Tokenize:}
 \emph{Mr.\ O'Neill
thinks that the boys' stories about Chile's capital aren't amusing.}
\end{frame}

\begin{frame}
\frametitle{Problemas de tokenização: Uma ou duas palavras? (ou várias)}
\begin{itemize}[<+->]
\item Hewlett-Packard
\item State-of-the-art
\item co-education
\item the hold-him-back-and-drag-him-away maneuver
\item data base
\item San Francisco
\item Los Angeles-based company
\item cheap San Francisco-Los Angeles fares
\item York University vs.\ New York University
\end{itemize}
\end{frame}


\begin{frame}
\frametitle{Números}
\begin{itemize}[<+->]
\item 3/20/91
\item 20/3/91
\item Mar 20, 1991
\item B-52
\item 100.2.86.144
\item (800) 234-2333
\item 800.234.2333
\item Sistemas mais antigos de RI costumavam não indexar números\ldots
\item \ldots  mas as vezes isto pode ser útil.
\end{itemize}
\end{frame}



\begin{frame}
\frametitle{Chinês: Ná há espaços}

\includegraphics{Chinese-example.eps}

\end{frame}

\begin{frame}
\frametitle{Segmentação ambígua em chinês}

\includegraphics[width=9cm]{ChineseAmbiguous.eps}

Os dois caracteres podem ser tratados como uma palavra significando `monge' 
ou como uma sequência de duas palavras significando ``e'' e ``parado''.


\end{frame}

\begin{frame}
\frametitle{Outros casos de ausência de espaços}
\begin{itemize}[<+->]
\item Palavras compostas em Holandês, Alemão, Sueco in Dutch, German, Swedish
\item Computerlinguistik $\rightarrow$ Computer + Linguistik
\item Lebensversicherungsgesellschaftsangestellter
\item $\rightarrow$ leben + versicherung + gesellschaft + angestellter
\item Inuit: tusaatsiarunnanngittualuujunga (Não consigo escutar direito.)
\item Muitas outras línguas apresentam dificuldades de segmentação: 
Finlandês, Urdu, \ldots
\end{itemize}
\end{frame}


\begin{frame}
\frametitle{Japonês}

\includegraphics[width=11cm]{Japanese-example-bitmap.eps}


\visible<2->{4 ``alfabetos'' diferentes: Caracteres chineses, 
hiragana silabário para finalizações infleccionais e palavras 
funcionais, 
katakana silabário para transcrições de palavras estrangeiras. Sem espaços 
(como no chinês).
}

\visible<3->{
Um usuário pode fazer uma consulta inteiramente em hiragana!
}

\end{frame}

\begin{frame}
\frametitle{Árabe}
\includegraphics[width=9cm]{Arabic-example.eps}
\end{frame}

\begin{frame}
\frametitle{Árabe: Bidirecionalidade}
\includegraphics[width=9cm]{Arabic-bidirectional.eps}

\visible<2->{
Bidirecionalidade não é um problema se o texto estiver codificado em 
Unicode.
}

\end{frame}

\begin{frame}
\frametitle{Acentos e diacríticos}
\begin{itemize}[<+->]
\item Acentos: Árvore vs.\ Arvore (Simples omissão de acento)
\item Tremas: Universit\"{a}t vs.\ Universitaet
  (substituição por sequencia especial de caracteres: ``ae'')
\item Critério mais importante: Qual a forma mais comumente adotada por 
usuários em suas consultas?
\item Mesmo em línguas que possuem acentos por padrão, freuentemente os 
usuários não os digitam. 
\end{itemize}
\end{frame}

\subsection{English}

\begin{frame}
\frametitle{Capitalização}
\begin{itemize}[<+->]
\item Reduzir todas as letras para minúsculas
\item Possíveis exceções: Palavras capitalizadas no meio da frase
\item MIT vs.\ mit
\item Fed vs.\ fed
\item Frequentemente é melhor converter tudo para minúsculas pois os 
usuários normalmente fazem suas consultas sem utilizar a capitalização 
correta.
\end{itemize}
\end{frame}

\begin{frame}
\frametitle{``Stop words''}
\begin{itemize}[<+->]
\item ``stop words'' = Palavras extremamente comuns que podem parecer ser 
 de pouco valor na seleção de documentos correspondentes à necessidade do 
usuário.
\item Exemplos: artigos, preposições, etc.

\item Eliminação de ``Stop word'' costumava ser padrão em sistemas de RI mais 
antigos.
\item Mas estas palavras são necessárias para consultas frasais,
  p.ex. ``Reino da Dinamarca''
\item A maioria dos engines de busca na web, indexa ``stop words''.
\end{itemize}
\end{frame}

\begin{frame}
\frametitle{Mais classes de equivalência}
\begin{itemize}[<+->]
\item Soundex: equivalencia fonética, Muller = Mueller)
\item Thesauri: equivalencia semântica, carro = automóvel)
\end{itemize}
\end{frame}


\begin{frame}
\frametitle{Lematização}
\begin{itemize}[<+->]
\item Reduz variações e inflexões a uma forma básica
\item Exemplo: \emph{sou, são, é $\rightarrow$ ser}
\item Exemplo: \emph{carro, carros $\rightarrow$ carros}
\item Exemple: \emph{Os carros do rapaz são de diferentes cores}
  $\rightarrow$ \emph{o carro do rapaz ser de diferente cor}
\item Lematização implica fazer a redução ``correta'' às forma mais fundamental 
segundo o dicionário {o \color{blue}lemma}).
\item Morfologia inflexional (\emph{cortando} 
$\rightarrow$
  \emph{cortar}) vs.\ morfologia derivacional (\emph{destruição} 
$\rightarrow$ \emph{destruir})
\end{itemize}
\end{frame}

\begin{frame}
\frametitle{Truncagem (``Stemming'')}
\begin{itemize}[<+->]
\item Definição de truncagem: Heurística bruta que {\color{blue}corta fora as 
terminações das palavras} na esperança de alcançar o mesmo objetivo da 
lematização sem a necessidade de conhecimento linguístico.
\item Depende da língua
\item Frequentemente inflexional {\color{blue}e} derivacional
\item Exemplo de truncagem derivacional: \emph{automatizar, automático,
  automatização} se reduzem a \emph{automat}
\end{itemize}
\end{frame}

\begin{frame}
\frametitle{Algoritmo de Porter}
\begin{itemize}[<+->]
\item Algoritmo mais comum para truncagem do inglês 
\item Resultados sugerem que é tão bom quanto outras opções de 
truncagem
\item Convenções + 5 fases de reduções 
\item Fases são aplicadas sequencialmente
\item Cada fase consiste de um conjunto de comandos.
\begin{itemize}[<+->]
\item Exemplo de comando: Deleta o \emph{ement} final se o que sobrar for mais 
longo do que um caractere
\item replacement $\rightarrow$ replac
\item cement $\rightarrow$ cement
\end{itemize}
\item Exemplo de convenção: Das regras em comando composto, selecione aquela 
que se aplica ao sufixo mais longo.
\end{itemize}
\end{frame}

\begin{frame}
\frametitle{Trucador Porter: Algumas regras}

\begin{tabular}[t]{lll@{\hspace{1in}}lll}
\multicolumn{3}{l}{\textbf{Regra}} & \multicolumn{3}{@{}l}{\textbf{Exemplo}} \\
SSES & $\rightarrow$ & SS  &   caresses  & $\rightarrow$ &  caress \\
IES  & $\rightarrow$ & I   &   ponies    & $\rightarrow$ &  poni \\
SS   & $\rightarrow$ & SS  &   caress    & $\rightarrow$ &  caress \\
S    & $\rightarrow$ &     &   cats      & $\rightarrow$ &  cat \\
\end{tabular}

\end{frame}

\begin{frame}[shrink=10]
\frametitle{Três Truncadores: Uma comparação}

\begin{description}

\item[\emph{Texto exemplo:}]
Such an analysis can reveal features that are not easily visible from
the variations in the individual genes and can lead to a picture of
expression that is more biologically transparent and accessible to
interpretation

\item[\emph{Truncador Porter:}]
such an analysi can reveal featur that ar not easili visibl from the
variat in the individu gene and can lead to a pictur of express that is
more biolog transpar and access to interpret

\item[\emph{Truncador Lovins:}]
such an analys can reve featur that ar not eas vis from
th vari in th individu gen and can lead to a pictur of
expres that is mor biolog transpar and acces to
interpres

\item[\emph{Truncador Paice:}]
such an analys can rev feat that are not easy vis from 
the vary in the individ gen and can lead to a pict of 
express that is mor biolog transp and access to 
interpret 

\end{description}

\end{frame}

\begin{frame}
\frametitle{Truncagem melhora efetividade?}
\begin{itemize}[<+->]
\item Em geral,  a truncagem aumenta a efetividade para algumas 
consultas  e diminui para outras.
\item Consultas em que a truncagem pode ajudar: [tartan
  sweaters], [sightseeing tour san francisco]
\item (classes de equivalência: \{sweater,sweaters\}, \{tour,tours\})
\item Classes de equivalência para \emph{oper} do truncador Porter 
 contém todas estas formas: \emph{
operate operating operates operation operative operatives operational}.
\item Consultas em que a truncagem prejudica: [operational AND research], 
[operating
  AND system], [operative AND dentistry]
\end{itemize}
\end{frame}


\begin{frame}
\frametitle{Exercise: O que o Google faz?}
\mygreen{\begin{itemize}
\item ``Stop words''
\item Normalização
\item Tokenização
\item minusculização
\item Truncagem
\item Alfabetos não latinos
\item Tremas
\item Palavras compostas
\item Números
\end{itemize}
}
\end{frame}
\section{Consultas de frases}

\begin{frame}
\frametitle{Consultas de frases}
\begin{itemize}[<+->]
\item Queremos responder a uma consulta como [stanford
  university] -- como uma frase.
\item Logo \emph{The inventor Stanford
  Ovshinsky never went to university} \myblue{não} deve ser uma resposta.
\item O conceito de consulta de frase é facilmente compreensível para 
usuários.
\item Cerca de 10\% das consultas na web são consultas de frases.
\item COnseguências para o índice invertido: Não é mais suficiente guardar 
docIDs  em listas de postings.
\item Duas maneiras de estender os índice invertido:
\begin{itemize}
\item índice de bigramas
\item Índice posicional
%\item \mygreen{Any ideas?}
\end{itemize}
\end{itemize}
\end{frame}

\begin{frame}
\frametitle{Índices de bigramas}
\begin{itemize}[<+->]
\item Indexa cada par consecutivo de termos no texto como uma frase.
\item Por exemplo, \emph{Friends, Romans, Countrymen} vai gerar dois 
bigramas: \emph{``friends romans''} e 
  \emph{``romans countrymen''}
\item Cada um destes bigramas é agora um termo do vocabulário.
\item Agora consultas de frases de duas palavras podem ser facilmente 
respondidas. 
\end{itemize}
\end{frame}


\begin{frame}
\frametitle{Consultas de frases mais longas}
\begin{itemize}[<+->]
\item Uma frase longa como \emph{``stanford university palo
  alto''} pode ser representada como a consulta Booleana 
\term{``stanford university''}
\oper{AND}
\term{``university palo''}
\oper{AND}
\term{``palo alto''}
\item É necessário pós-filtrar o conjunto de respostas para identificar o 
subconjunto que contém a frase de quatro palavras.
\end{itemize}
\end{frame}

\begin{frame}
\frametitle{Problemas com índices de bigramas}
\begin{itemize}[<+->]
\item \mygreen{Porque estes índices são raramente usados?}
\item Falsos Positivos, como indicado acima
\item Índice explode devido ao tamanho exagerado do vocabulário
\end{itemize}
\end{frame}

\begin{frame}[<+->]
\frametitle{Índices posicionais}
\pause[2]
\begin{itemize}
\item Índices posicionais são mais eficientes que indices de bigramas.
\item Listas de Postings em um índice {\color{blue}não posicional}: Cada 
posting é apenas um docID
\item Listas de Postings em um índice {\color{blue} posicional}: Cada 
posting é um docID e {\color{blue} uma lista de posições}

\end{itemize}

\end{frame}

\begin{frame}
\frametitle{Índices positionais: Exemplo}

\pause[2]

Query: \emph{``to$_1$ be$_2$ or$_3$ not$_4$ to$_5$ be$_6$''}

\pause[3]

\raggedright
\term{to}, 993427:\\
\hspace*{2em}\llap{$\langle$} \alert<2-6>{1}: $\langle$\alert<2,3>{7}, 
\alert<4,5>{18}, \alert<6>{33, 72, 86, 231}$\rangle$;\\
\hspace*{2em} \alert<7>{2}: $\langle$1, 17, 74, 222, 255$\rangle$; \\
\hspace*{2em} \alert<8->{4}: $\langle$\alert<9>{8}, \alert<10-12>{16}, 
\alert<12,13,14,15>{190}, \alert<15->{429}, \alert<17->{433}$\rangle$;\\
\hspace*{2em} 5: $\langle$363, 367$\rangle$; \\
\hspace*{2em} 7: $\langle$13, 23, 191$\rangle$; 
\ldots $\rangle$ \\[2ex]
\term{be}, 178239:\\
\hspace*{2em}\llap{$\langle$}  \alert<2-6>{1}: $\langle$\alert<3>{17}, 
\alert<5>{25}$\rangle$;\\
\hspace*{2em} \alert<8->{4}: $\langle$\alert<11-12>{17}, \alert<14,15>{191}, 
291, \alert<16->{430}, \alert<18->{434}$\rangle$;\\
\hspace*{2em} 5: $\langle$14, 19, 101$\rangle$; 
\ldots $\rangle$



\smallskip

\visible<19->{Documento 4 é a resposta!}

\end{frame}

\begin{frame}
\frametitle{Busca por proximidade }
\begin{itemize}[<+->]
\item Acabamos de ver como usar um índice posicional para busca de 
frases.
\item Também podemos usá-los para busca por proximidade.
\item Por exemplo: \query{employment /4 place}
\item Encontrar todos os documentos que contêm \term{employment} e 
  \term{place} a até 4 palavras words de distância entre si.
\item \emph{Employment agencies that place healthcare
  workers are seeing growth} é uma resposta.
\item \emph{Employment agencies that have learned to adapt
  now place healthcare
  workers} não é uma resposta.
\end{itemize}
\end{frame}

\begin{frame}
\frametitle{Busca por proximidade}
\begin{itemize}[<+->]
\item Use o índice posicional
\item Algoritmo mais simples: Olhe para o Produto vetorial das 
posições de (i) \term{employment} no documento e (ii)
  \term{place} no documento
\item Muito ineficiente para palavras frequentes, especialmente ``stop 
words''
\item Note que queremos retornar as posições encontradas e não apenas 
uma lista de documentos.
\item Isto é importante para sumarizações dinâmicas etc.
\end{itemize}
\end{frame}

\begin{frame}[shrink=20]
\frametitle{Interseção de ``proximidade'' }


\begin{algorithm}{PositionalIntersect}{p_1, p_2, k}
answer \= \langle\;\rangle \\
\begin{WHILE} {p_1 \neq \NIL \text{and} p_2 \neq \NIL}
  \begin{IF}{docID(p_1) = docID(p_2)}
    l \= \langle\;\rangle\\
    pp_1 \= positions(p_1) \\
    pp_2 \= positions(p_2) \\
    \begin{WHILE} {pp_1 \neq \NIL}
      \begin{WHILE} {pp_2 \neq \NIL}
        \begin{IF} {|pos(pp_1)-pos(pp_2)| \le k}
          \CALL{Add}(l,pos(pp_2))
        \ELSE \begin{IF} {pos(pp_2) 
%%<
> pos(pp_1)}
                 \algkey{break}
              \end{IF}
        \end{IF} \\
        pp_2 \= next(pp_2)
      \end{WHILE} \\
      \begin{WHILE} {l \neq \langle\;\rangle \text{and}
          |l[0]-pos(pp_1)|
%%<
>k}
        \CALL{Delete}(l[0])
      \end{WHILE} \\
      \begin{FOR} {\EACH ps \in l}
        \CALL{Add}(answer, \langle docID(p_1), pos(pp_1), ps\rangle) 
      \end{FOR}\\
      pp_1 \= next(pp_1)
    \end{WHILE} \\
    p_1 \= next(p_1) \\
    p_2 \= next(p_2)
  \ELSE
    \begin{IF} {docID(p_1) < 
%%>
docID(p_2)}
      p_1 \= next(p_1)
    \ELSE 
      p_2 \= next(p_2)
    \end{IF}
  \end{IF}
\end{WHILE} \\
\RETURN{answer}
\end{algorithm}

\end{frame}

\begin{frame}
\frametitle{Esquema Combinado}
\begin{itemize}[<+->]
\item Índices de bigramas e posicionais podem ser combinados proveitosamente.
\item Muitos bigramas são extremamente frequentes: Michael Jackson,
  Britney Spears,  etc.
\item Para estes bigramas, o aumento de velocidade proporcionado pela 
indexação de bigramas, é substancial.
\item Esquema de combinação: Incluir bigramas frequentescomo 
vocabulário no índice. Todas as outras frases são recuperadas por interseção 
posicional.
\end{itemize}
\end{frame}

\begin{frame}
\frametitle{Consultas ``Posicionais'' no Google}
\begin{itemize}[<+->]
\item Para buscas na Web, consultas posicionais são muito mais caras do que 
consultas Booleanas regulares.
\item Vejamos o exemplo de consultas de frases.
\item {\color{green} Porque são mais custosas?}
\item {\color{green}Você consegue demonstrar no Google que as consultas de 
frases são mais custosas que consultas Booleanas?}
\end{itemize}
\end{frame}


\againframe<6>{takeaway}

\end{document}
